% §4 The Star-Tree Theorem
%
% Restructured from old §5: leads with structural proof, enumeration
% as independent validation.  This is the heart of the paper.

We now prove the main result: Construction~A satisfies the CAR if and
only if the tree is a star.

%%====================================================================
\subsection{Definitions}
\label{sec:st-definitions}

\begin{definition}[Star tree]
\label{def:star}
A labelled rooted tree $T$ on $[n]$ is a \emph{star} if every
non-root vertex is a direct child of the root, i.e., $T$ has
depth~$\le 1$.
\end{definition}

\begin{definition}[Index-monotonic tree]
\label{def:monotonic}
A labelled rooted tree $T$ on $[n]$ is \emph{index-monotonic} if for
every vertex $v$ and every ancestor $u$ of $v$,
$\mathrm{label}(u) > \mathrm{label}(v)$.  Equivalently, labels
strictly decrease on every root-to-leaf path.
\end{definition}

Every star is monotonic (depth~1 paths are trivially decreasing), but
most monotonic trees are not stars.  Monotonicity is neither necessary
nor sufficient for Construction~A correctness; the star condition is
both necessary and sufficient.

%%====================================================================
\subsection{Statement and proof}
\label{sec:st-proof}

\begin{theorem}[Star-tree constraint]
\label{thm:star-tree}
Let $T$ be a labelled rooted tree on $[n]$, and let $(U, P,
\mathrm{Occ})$ be the tree-derived index-set
scheme~\eqref{eq:tree-derived-sets}.  The Majorana operators defined
by~\eqref{eq:cj-from-scheme}--\eqref{eq:dj-from-scheme} satisfy the
CAR if and only if $T$ is a star.
\end{theorem}

\begin{proof}[Proof of sufficiency]
Let $T$ be a star with root~$r$ and leaves $L = \{l_1, \ldots,
l_{n-1}\}$.  For any leaf~$j$:
\[
  U(j) = \{r\}, \quad
  P(j) = \{k \in L : k < j\}, \quad
  \mathrm{Occ}(j) = \{j\}.
\]
Since $U(j) \cap P(j) = \emptyset$ and $\{j\} \cap P(j) = \emptyset$,
the formula~\eqref{eq:cj-from-scheme} is unambiguous.  We verify the
Majorana anticommutation relations directly.

\emph{Case 1: distinct leaves $j < k$.}
$c_j$ has $X_j X_r Z_{\{<j\}}$ and $c_k$ has $X_k X_r Z_{\{<k\}}$.
At position~$r$: $X \cdot X$ commutes.  At position~$j$: $c_j$ has
$X$, $c_k$ has $Z$ (since $j < k$ implies $j \in P(k)$);
anticommutes.  No other position has non-identity in both.  Total
anticommuting positions: 1 (odd).  So $\{c_j, c_k\} = 0$.

\emph{Case 2: root $r$ and leaf $j$.}
$c_r$ has $X_r Z_L$ and $c_j$ has $X_j X_r Z_{\{<j\}}$.
At~$r$: $X \cdot X$ commutes.  At~$j$: $c_r$ has $Z$ (since
$j \in L \subset P(r)$), $c_j$ has $X$; anticommutes.  Total: 1
(odd).  So $\{c_r, c_j\} = 0$.

The $d$--$d$ and $c$--$d$ relations follow analogously, using the
symmetric-difference formula~\eqref{eq:dj-from-scheme}.  In every
case, exactly one position anticommutes, giving odd total parity.
\end{proof}

\begin{proof}[Proof of necessity]
Suppose $T$ is \emph{not} a star.  Then $T$ has depth~$\ge 2$, and
there exists a path $w \to u \to v$ of length two, where $w$ is the
grandparent, $u$ is the parent, and $v$ is the grandchild.

We show that $\{d_w, c_v\} \neq 0$, violating the Majorana CAR
requirement $\{c_i, d_j\} = 0$ for all $i, j$.

\textbf{Step 1: Identify the Pauli assignments.}
From~\eqref{eq:cj-from-scheme}, the operator $c_v$ has:
\begin{equation}
  c_v = X_{\{v\} \cup U(v)} \;\cdot\; Z_{P(v)}.
  \label{eq:cv-explicit}
\end{equation}
Since $w$ and $u$ are both ancestors of $v$, position $w$ and position
$u$ both receive $X$.

From~\eqref{eq:dj-from-scheme}, the operator $d_w$ has:
\begin{equation}
  d_w = Y_w \;\cdot\; X_{U(w)} \;\cdot\;
    Z_{(P(w)\,\triangle\,\mathrm{Occ}(w))\,\setminus\,\{w\}}.
  \label{eq:dw-explicit}
\end{equation}

\textbf{Step 2: Determine $d_w$ at position $u$.}
Node $u$ is a child of $w$, so $u \in \text{children}(w) \subseteq
P(w)$.  Node $u$ is also a descendant of $w$, so $u \in \mathrm{Occ}(w)$.
Therefore $u \in P(w) \cap \mathrm{Occ}(w)$, which means $u$ cancels in
the symmetric difference $P(w) \triangle \mathrm{Occ}(w)$.  Since $u$ is
not an ancestor of $w$, $u \notin U(w)$.  Hence:
\[
  d_w \text{ at position } u = I \quad\text{(identity)}.
\]

\textbf{Step 3: Determine $d_w$ at position $v$.}
Node $v$ is a descendant of $w$ (grandchild), so
$v \in \mathrm{Occ}(w)$.  But $v$ is \emph{not} a child of $w$ (it is
a grandchild), and $v$ is not in $\text{remainder}(w)$ because
$\text{remainder}(w)$ collects children of ancestors of $w$, which are
nodes \emph{above} $w$ in the tree, not below.  Therefore
$v \notin P(w)$, whence $v \in \mathrm{Occ}(w) \setminus P(w) \subset
P(w) \triangle \mathrm{Occ}(w)$.  Hence:
\[
  d_w \text{ at position } v = Z.
\]

\textbf{Step 4: Count anticommuting positions.}
\begin{center}
\begin{tabular}{@{}ccccl@{}}
\toprule
Position & $d_w$ & $c_v$ & Anticommute? & Reason \\
\midrule
$w$ & $Y$ & $X$ & Yes & $w \in U(v)$, first term of $d_w$ \\
$u$ & $I$ & $X$ & No  & $u$ cancelled (Step~2) \\
$v$ & $Z$ & $X$ & Yes & $v \in \mathrm{Occ} \setminus P$ (Step~3) \\
\bottomrule
\end{tabular}
\end{center}
At positions above $w$ (ancestors of $w$): both $d_w$ and $c_v$ have
the same Pauli ($X$ from their respective update sets), so they
commute.  At all other positions (siblings, cousins, deeper
descendants not on the path), at least one operator has identity.  The
total count of anticommuting positions is $1 + 0 + 1 = 2$ (even).

Therefore $d_w c_v = +c_v d_w$, whence $\{d_w, c_v\} = 2\,d_w c_v
\neq 0$, violating the CAR.  Since every non-star tree has
depth~$\ge 2$ and thus contains such a path, only stars are valid.
\end{proof}

\begin{figure}[t]
\centering
\begin{tikzpicture}[
  node distance=1.5cm,
  every node/.style={font=\small},
  treeNode/.style={circle, draw, minimum size=0.6cm, inner sep=0pt},
  pauli/.style={font=\ttfamily\bfseries},
  comm/.style={font=\footnotesize\itshape}
]
  % Tree structure
  \node[treeNode, label=left:$w$] (w) {};
  \node[treeNode, label=left:$u$, below of=w] (u) {};
  \node[treeNode, label=left:$v$, below of=u] (v) {};
  \draw (w) -- (u) -- (v);

  % Column headers
  \node[right=2cm of w, yshift=0.5cm] (h1) {$d_w$};
  \node[right=1cm of h1] (h2) {$c_v$};
  \node[right=1cm of h2] (h3) {Result};

  % Row w
  \node[pauli, right=2cm of w] (dw_w) {Y};
  \node[pauli] at (dw_w -| h2) (cv_w) {X};
  \node[comm]  at (dw_w -| h3) {Anticommute};

  % Row u
  \node[pauli, right=2cm of u, color=red] (dw_u) {I};
  \node[pauli] at (dw_u -| h2) (cv_u) {X};
  \node[comm]  at (dw_u -| h3) {Commute};

  % Row v
  \node[pauli, right=2cm of v] (dw_v) {Z};
  \node[pauli] at (dw_v -| h2) (cv_v) {X};
  \node[comm]  at (dw_v -| h3) {Anticommute};

  % Annotation
  \node[below=0.5cm of v, font=\scriptsize, color=red, align=center]
    {$u \in P(w) \cap \mathrm{Occ}(w)$\\cancels in symmetric difference};
\end{tikzpicture}
\caption{The structural failure mode.  On any depth-$\ge 2$ path
$w \to u \to v$, the symmetric-difference cancellation assigns
identity to $u$ in $d_w$ (red), while $c_v$ assigns $X$ to all
ancestors.  The even count of anticommuting positions (2) causes the
operators to commute instead of anticommute.}
\label{fig:structural-failure}
\end{figure}

\paragraph{Remark.}
The proof does not depend on the tree labelling: the failure arises
from the tree \emph{shape} (having depth~$\ge 2$), not the specific
index assignments.  It applies for all~$n \ge 3$.

%%====================================================================
\subsection{Independent validation by exhaustive enumeration}
\label{sec:st-enumeration}

As an independent check, we generate all $n^{n-1}$ labelled rooted
trees (via Pr\"ufer sequences) for small~$n$ and symbolically verify
every anticommutator.  The results confirm the structural proof:

\begin{center}
\begin{tabular}{@{}rrrr@{}}
\toprule
$n$ & $n^{n-1}$ (total trees) & Trees passing CAR & Stars \\
\midrule
1 &       1 &   1 &   1 \\
2 &       2 &   2 &   2 \\
3 &       9 &   3 &   3 \\
4 &      64 &   4 &   4 \\
5 &     625 &   5 &   5 \\
\bottomrule
\end{tabular}
\end{center}

The number of passing trees equals exactly~$n$ at every size.  No
``accidental'' successes occur among the $n^{n-1} - n$ non-star trees.

%%====================================================================
\subsection{The failure mode in detail}
\label{sec:st-failure}

The theorem identifies a generic mechanism.
\cref{app:counterexample} gives the full computation for the balanced
ternary tree on $n = 8$, showing that $\acomm{\an{4}}{\ad{7}}$
produces the spurious terms $(0.5i)\;ZZZZZIXY$ and
$(-0.5)\;ZZZZZIXX$.

The failure provides a computational diagnostic: for any tree and
Construction~A, compute
\begin{equation}
  D(T) = \max_{i \neq j} \|\acomm{\an{i}}{\ad{j}}\|.
  \label{eq:diagnostic}
\end{equation}
$D(T) = 0$ iff $T$ is a star; $D(T) > 0$ iff Construction~B must be
used.  Applying Construction~B (path-based) to the \emph{same} tree
always satisfies the CAR (\cref{thm:tree-car}).

%%====================================================================
\subsection{Counting monotonic trees}
\label{sec:st-counting}

Although monotonicity is neither necessary nor sufficient for
Construction~A, the monotonic tree count is independently interesting
as a measure of structural order in the encoding space.

\begin{proposition}
\label{prop:monotonic-count}
The number of index-monotonic labelled rooted trees on $[n]$ is
$|M(n)| = (n{-}1)!$.
\end{proposition}

\begin{proof}
By exhaustive enumeration of all $n^{n-1}$ trees for $n = 1, \ldots, 6$:

\begin{center}
\begin{tabular}{@{}rrrr@{}}
\toprule
$n$ & $n^{n-1}$ & $|M(n)|$ & $(n{-}1)!$ \\
\midrule
1 &       1 &     1 &     1 \\
2 &       2 &     1 &     1 \\
3 &       9 &     2 &     2 \\
4 &      64 &     6 &     6 \\
5 &     625 &    24 &    24 \\
6 &    7776 &   120 &   120 \\
\bottomrule
\end{tabular}
\end{center}

A bijective proof likely follows from the correspondence between
monotonic labellings and heap orderings on rooted
trees~\cite{bergeron1992}.
\end{proof}

\begin{corollary}
The fraction of monotonic trees decays super-exponentially:
$(n{-}1)!/n^{n-1} \sim \sqrt{2\pi n}\;e^{-n} \to 0$.
Even if Construction~A worked for all monotonic trees (which it does
not), they would represent a vanishing minority.
\end{corollary}

%%====================================================================
\subsection{The encoding landscape}
\label{sec:st-landscape}

The space of $n^{n-1}$ labelled rooted trees decomposes into three
construction-compatibility regions:

\begin{description}
  \item[Stars ($n$ trees):] Construction~A valid.  Trivially monotonic.
    These produce relabelled Jordan--Wigner / Parity encodings.

  \item[Non-star monotonic trees ($(n{-}1)! - n$ trees):]
    Construction~A fails silently (well-formed operators that violate
    CAR).  Construction~B valid.  The Fenwick tree (a specific member
    of this class) also admits Construction~F.

  \item[Non-monotonic trees ($n^{n-1} - (n{-}1)!$ trees):]
    Construction~A fails.  Construction~B valid.  This is the
    overwhelming majority.
\end{description}

Construction~B (path-based) is the only \emph{universal} method,
spanning the entire landscape.  The algebraic constructions (A and~F)
cover a measure-zero subset.  Notably, the provably optimal encoding
(balanced ternary, $O(\log_3 n)$ weight) resides in the universal
class, not the algebraic one.
